\subsection{Fuzzy Anonymous Complaint Tally System (FACTS)}\label{sec:facts}

In this section, we present the syntax for FACTS and describe how FACTS is used.  We show how to instantiate FACTS in Section~\ref{sec:construction}.

\paragraph{Assumptions:}
We assume that each user $A$ has a unique identifier $ID_A$, and that the server can authenticate these IDs.  (We will abuse notation to use $A$ to represent the user and also the id $ID_A$).  We also assume that the server has an identifier $ID_S$ (we will denote this by $S$) that can be authenticated by all users.

Additionally, we assume that the underlying end-to-end encrypted messaging system (EEMS) offers methods $\sendee(A,B,x)$ and $\receiveee(A,B,x)$ for sending and verifying a message $x$ sent from user $A$ to user $B$.  Moreover, we assume that this communication is encrypted and authenticated. In particular, $\receiveee$ verifies that the received message was sent by $A$ and was not modified in transit.  Importantly, we do not assume that this platform is anonymous, instead assuming that the full messaging history i.e., who sent a message to whom and the size of that message, is available to the server. 

\paragraph{Syntax:}
FACTS is a tuple of protocols $\FACTS=(\setup, \send, \receive, \complain, \audit)$.  The first is used to set up FACTS, the next two are used to send and verify messages, while the last two methods are used to issue complaints and audit received messages.

\begin{itemize}
    \item $\setup(c)$:  This takes as input an upper bound on the total number of users and initiates the FACTS scheme for $c$ users.
    
    \item $\send(A, B,\tag_x,x)$: This method is used by a user $A$ to send a message $x$ to another user $B$.  This may be a new message \emph{originated} by $A$ (indicated by $\tag_x = \perp$) or a forward of a previously received message.  
    
    
    \item $\receive(A,B,\tag_x, x)$:  This algorithm is run by $B$ upon receiving a message $(\tag_x, x)$ from $A$.
    
    This algorithm checks whether $\tag_x$ is indeed a valid tag generated by $A$ on message $x$.  If this is the case, then $B$ accepts the message, otherwise he rejects the message.
    
    \item $\complain(C, \tag_x, x)$:  This protocol is run by a user $C$ to complain about a received message $(\tag_x,x)$.  
    
    \item $\audit(C,\tag_x, x)$:  This protocol issues an audit of a message $x$ revealing $(\tag_x, x)$ to $S$.  This will be called by $C$ when the number of complaints on $m$ exceeds a pre-defined threshold (with high probability).
\end{itemize}

\paragraph{Usage:}
%\anote{Not sure we really need this paragraph.}
The following workflow demonstrates the standard usage of FACTS.  To originate a new message $x$, a user $A$ runs the $\send$ protocol with the server $S$ to create metadata $\tag_x$.  $\send$ then sends this metadata and the message $(\tag_x,x)$
to the receiving user $B$ using the messaging platforms $\sendee$ method.  Upon receiving a message $(\tag_x, x)$, $B$ first locally runs $\receive(A,B,\tag_x,x)$ to verify that the received message and tag are valid, if this fails he ignores the message.  To forward a received message $(\tag_x,x)$, a user $A$ runs $\send$ with the server $S$ to produce metadata $\tag'_x$; $A$ then discards this metadata, and the original message $(\tag_x, x)$ is sent instead using the messaging platform's $\sendee$ method.\footnote{We note that since the underlying messaging scheme is encrypted, the actual ciphertext sent will not be the same as the ciphertext received.}

If a user $B$ receives a message $(\tag_x, x)$ that it considers ``fake'', he can use the  $\complain$ protocol to issue a new complaint on this message.  After issuing a complaint, $B$ checks whether the threshold of complaints on $x$ has been reached.  If so, he calls $\audit$ to trigger an audit on the message $(\tag_x,x)$, revealing $x$ and the originator of $x$ to the server $S$.

We note that users may join and leave during the execution of FACTS as long as the total number of identifiable users does not exceed $c$.

% \subsubsection{Security of FACTS}
% We now provide definitions for the security achieves by FACTS.  As discussed previously, we aim to achieve the following security goals.  
% \begin{itemize}
%     \item Privacy of messages and complaints should guarantee that the server $S$ together with a set of colluding clients should not learn anything about the contents of any message they have not received. Additionally, privacy of complaints should guarantee that the server (and colluding clients) should not learn any information about the pattern of complaints until an audit on a particular message is triggered.
%     \item Integrity of originator guarantees that upon an audit only the originator of the message or another malicious party will be identified as an originator
%     \item Finally, integrity of complaints guarantees that a malicious client who has not received $m$ cannot significantly speed up or slow down complaints for $m$.
% \end{itemize}
